\documentclass[11pt]{article}
\usepackage{amsmath,amssymb,amsthm}
\usepackage[margin=1in]{geometry}

\newtheorem{theorem}{Theorem}
\newtheorem{lemma}[theorem]{Lemma}
\newtheorem{definition}[theorem]{Definition}
\newtheorem{proposition}[theorem]{Proposition}
\newtheorem{remark}[theorem]{Remark}

\title{Internal Verification of Proof Checking\\
in Guard--Rose Recursive Arithmetic:\\
Experiments Toward Nelson's Program}
\author{Summary of Agda development}
\date{April 2026}

\begin{document}
\maketitle

\section{Introduction}

We investigate, within a formalization of G\"odel's second incompleteness
theorem for recursive arithmetic, the extent to which the system can
\emph{internally verify its own proof-checking function}. This is
motivated by Nelson's program~\cite{nelson}, which asks whether the
consistency of a system such as primitive recursive arithmetic can be
proved within the system itself---and whether the resulting self-application
leads to a contradiction.

The formalization uses an equational system of binary-tree recursive
arithmetic that combines features of Rose~\cite{rose1961} and
Guard~\cite{guard1963}. From Rose we take the tree-based data domain,
the coding scheme, and the equational approach to G\"odel's theorems.
From Guard we take the formulation with 1-ary and 2-ary functors
(following Church's Binary Recursive Arithmetic) and the explicit
Deriv-level proof system. The system is formalized in Agda with
\texttt{--safe --without-K --exact-split}. No postulates are used.

\section{The System}

\subsection{Terms and Functors}

The language has:
\begin{itemize}
\item \textbf{Terms}: $O$ (leaf), $\mathit{var}\;n$, $f(t)$ for $f$ a
  1-functor, $g(t_1, t_2)$ for $g$ a 2-functor.
\item \textbf{1-functors} (\texttt{Fun1}): $I$, $\mathit{Fst}$,
  $\mathit{Snd}$, $\mathit{Comp}\;f\;g$, $\mathit{Comp2}\;h\;f\;g$,
  $\mathit{Rec}\;z\;s$, $\mathit{KT}\;t$.
\item \textbf{2-functors} (\texttt{Fun2}): $\mathit{Pair}$,
  $\mathit{Const}$, $\mathit{Lift}\;f$, $\mathit{Post}\;f\;h$,
  $\mathit{Fan}\;h_1\;h_2\;h$, $\mathit{IfLf}$, $\mathit{TreeEq}$.
\item \textbf{Equations}: $t = u$ for terms $t$, $u$.
\end{itemize}

Binary trees ($\mathit{lf}$ and $\mathit{nd}\;a\;b$) serve as the data
domain. $\mathit{reify} : \mathit{Tree} \to \mathit{Term}$ maps
$\mathit{lf} \mapsto O$ and $\mathit{nd}\;a\;b \mapsto
\mathit{Pair}(\mathit{reify}\;a, \mathit{reify}\;b)$.

The 1-functor $\mathit{KT}\;t$ is the constant function returning~$t$.
$\mathit{Rec}\;z\;s$ defines tree recursion with base value~$z$ and
step function~$s$.

\subsection{Derivation Rules}

$\mathit{Deriv}\;\mathit{hyp}\;\mathit{eq}$ is the type of derivations
of equation $\mathit{eq}$ from hypothesis $\mathit{hyp}$. The system has
26~constructors, organized as follows.

\medskip
\noindent\textbf{Computation axioms} (one-step reductions):
\begin{align}
  \mathit{axI}&: I(t) = t \tag{tag 0}\\
  \mathit{axFst}&: \mathit{Fst}(\mathit{Pair}(a,b)) = a \tag{tag 1}\\
  \mathit{axSnd}&: \mathit{Snd}(\mathit{Pair}(a,b)) = b \tag{tag 2}\\
  \mathit{axConst}&: \mathit{Const}(a,b) = a \tag{tag 3}\\
  \mathit{axComp}&: \mathit{Comp}(f,g)(t) = f(g(t)) \tag{tag 4}\\
  \mathit{axComp2}&: \mathit{Comp2}(h,f,g)(t) = h(f(t), g(t)) \tag{tag 5}\\
  \mathit{axLift}&: \mathit{Lift}(f)(a,b) = f(a) \tag{tag 6}\\
  \mathit{axPost}&: \mathit{Post}(f,h)(a,b) = f(h(a,b)) \tag{tag 7}\\
  \mathit{axFan}&: \mathit{Fan}(h_1,h_2,h)(a,b) = h(h_1(a,b), h_2(a,b)) \tag{tag 8}\\
  \mathit{axRecLf}&: \mathit{Rec}(z,s)(O) = z \tag{tag 9}\\
  \mathit{axRecNd}&: \mathit{Rec}(z,s)(\mathit{Pair}(a,b)) = s(\mathit{Pair}(a,b),
    \mathit{Pair}(\mathit{Rec}(z,s)(a), \mathit{Rec}(z,s)(b))) \tag{tag 10}
\end{align}

\noindent\textbf{Conditional and comparison axioms}:
\begin{align}
  \mathit{axIfLfL}&: \mathit{IfLf}(O, \mathit{Pair}(a,b)) = a \tag{tag 11}\\
  \mathit{axIfLfN}&: \mathit{IfLf}(\mathit{Pair}(x,y), \mathit{Pair}(a,b)) = b \tag{tag 12}\\
  \mathit{axTreeEqLL}&: \mathit{TreeEq}(O, O) = O \tag{tag 13}\\
  \mathit{axTreeEqLN}&: \mathit{TreeEq}(O, \mathit{Pair}(a,b)) = \mathit{Pair}(O,O) \tag{tag 14}\\
  \mathit{axTreeEqNL}&: \mathit{TreeEq}(\mathit{Pair}(a,b), O) = \mathit{Pair}(O,O) \tag{tag 15}\\
  \mathit{axTreeEqNN}&: \mathit{TreeEq}(\mathit{Pair}(a_1,a_2), \mathit{Pair}(b_1,b_2)) \notag\\
    &= \mathit{IfLf}(\mathit{TreeEq}(a_1,b_1),
       \mathit{Pair}(\mathit{TreeEq}(a_2,b_2), \mathit{Pair}(O,O))) \tag{tag 16}
\end{align}

\noindent\textbf{Structural rules}:
\begin{align}
  \mathit{axRefl}&: t = t \tag{tag 17}\\
  \mathit{ruleSym}&: t = u \;\Rightarrow\; u = t \tag{tag 18}\\
  \mathit{ruleTrans}&: t = u,\; u = v \;\Rightarrow\; t = v \tag{tag 19}\\
  \mathit{cong1}&: t = u \;\Rightarrow\; f(t) = f(u) \tag{tag 20}\\
  \mathit{congL}&: t = u \;\Rightarrow\; g(t,x) = g(u,x) \tag{tag 21}\\
  \mathit{congR}&: t = u \;\Rightarrow\; g(x,t) = g(x,u) \tag{tag 22}
\end{align}

\noindent\textbf{Substitution and hypothesis}:
\begin{align}
  \mathit{ruleInst}&: \mathit{Deriv}\;\mathit{hyp}\;\mathit{eq}
    \;\Rightarrow\; \mathit{Deriv}\;\mathit{hyp}\;(\mathit{eq}[x := t])
    \tag{tag 23}\\
  \mathit{ruleHyp}&: \mathit{Deriv}\;\mathit{hyp}\;\mathit{hyp}
    \tag{--}
\end{align}

\noindent\textbf{Schema~F} (uniqueness of tree recursion):
\begin{align}
  \mathit{ruleF}&: \text{if } f \text{ and } g \text{ both satisfy }
    \mathit{Rec}\;z\;s, \text{ then } f(x) = g(x) \tag{tag 24}
\end{align}
Schema~F has four premises: $f(O) = z$, $f(\mathit{Pair}(x,y)) =
s(\mathit{Pair}(x,y), \mathit{Pair}(f(x), f(y)))$, and the corresponding
two premises for~$g$.

\noindent\textbf{Constant function axiom}:
\begin{align}
  \mathit{axKT}&: \mathit{KT}(t)(x) = t \tag{tag 25}
\end{align}

\begin{remark}
Tags 0--25 correspond to the 26 constructors of $\mathit{Deriv}$ and are
used in proof-encoding trees processed by the internal proof
checker~$\mathit{thFunTFor}$. The tag assignment follows Guard's
approach~\cite{guard1963} of assigning G\"odel numbers to inference rules.
\end{remark}

\subsection{Relationship to Rose and Guard}

Rose~\cite{rose1961} works with Ternary Recursive Arithmetic (1-ary, 2-ary,
and 3-ary functors) over natural numbers. Guard~\cite{guard1963} works
with Binary Recursive Arithmetic (1-ary and 2-ary functors only), following
Church. Our system uses Guard's functor vocabulary but Rose's tree-based
data domain. The key correspondences are:

\begin{center}
\begin{tabular}{lll}
\hline
\textbf{Component} & \textbf{Rose/Guard} & \textbf{Ours} \\
\hline
Data domain & Natural numbers & Binary trees \\
Functors & 1-ary, 2-ary (Guard) & 1-ary, 2-ary \\
Recursion & Primitive recursion on $\mathbb{N}$ & Tree recursion (Rec) \\
Induction & Schema F (Rose) / simple induction (Guard) & Schema F \\
Equality test & $a(x,y) = 0$ iff $x = y$ & $\mathit{TreeEq}(t,t) = O$ \\
Proof evaluator & $th(y)$ (Rose) / $th(y)$ (Guard) & $\mathit{thFunTFor}$ \\
Coded substitution & $se(x,y,z)$ / $\mathit{sub}(z,P)$ & $\mathit{substTFor}$ \\
\hline
\end{tabular}
\end{center}

Guard~\cite{guard1963} proves both G\"odel incompleteness theorems for BRA.
His proof of the second theorem (Theorem~14 in Guard) formalizes the first
theorem's proof inside the system by showing that the proof predicate
$th(y)$ tracks through all inference rules. Our $\mathit{thm14E}$ is the
Agda formalization of this step.

\section{Internal Proof Checker}

\subsection{thFunTFor}

The internal proof checker $\mathit{thFunTFor} = \mathit{Rec}\;O\;\mathit{thFunStep}$
is a tree-recursive function defined entirely within the object language.
It maps proof-encoding trees to coded equations: given a tree encoding a
derivation, $\mathit{thFunTFor}$ returns $\mathit{Pair}(L, R)$ where
$L$ and $R$ are the codes of the left and right sides of the derived equation.

The step function $\mathit{thFunStep}$ is a 26-branch dispatch on the
encoding tag. Each branch extracts the relevant components from the
encoding and constructs the output equation code.

\subsection{Proof Encodings}

Each constructor of $\mathit{Deriv}$ has an encoding function
$\mathit{enc}_k : \ldots \to \mathit{Tree}$ that produces a tree of the
form $\mathit{nd}(\mathit{natCode}\;k)(\mathit{data})$, where $k$ is
the tag (0--25) and $\mathit{data}$ contains the sub-proof encodings
and auxiliary data. The encodings are strictly compositional: no
sub-encoding is duplicated.

\subsection{Derivation Size}

We define $\mathit{derivSize}$ on $\mathit{Deriv}$ and prove:

\begin{theorem}[DerivSize]
The Schema~F proof $\mathit{treeEqSelfAll}$ (which proves
$\mathit{TreeEq}(t,t) = O$ for all~$t$) uses exactly $113$ inference
steps. The instantiation $\mathit{treeEqSelf}\;t$ uses $125$ steps,
independent of~$t$.
\end{theorem}

Both results are verified by \texttt{refl} in Agda, confirming that the
Schema~F proof is constant-size.

\section{Nelson Experiments: Internal Verification}

Motivated by Nelson's program~\cite{nelson}, we investigate whether the
system can internally verify the correctness of its own proof checker.
Concretely, we prove \emph{inside} $\mathit{Deriv}$ that
$\mathit{thFunTFor}$ correctly handles specific encoding tags.

\subsection{Infrastructure}

\begin{definition}[Tag dispatch]
For each tag $k \in \{0, \ldots, 25\}$, the $\mathit{ndDispatch}$ chain
checks $\mathit{TreeEq}(\mathit{Fst}(\mathit{orig}),
\mathit{reify}(\mathit{natCode}\;k))$ against each tag in sequence.
A \emph{branch miss} (tag mismatch) falls through to the next branch;
a \emph{branch hit} (tag match) selects the corresponding case function.
\end{definition}

\begin{definition}[Passthrough]
When the left child of a tree node is $\mathit{Pair}$-tagged (i.e.,
$\mathit{Pair}(\mathit{Pair}(a_1,a_2), b)$), all 26 tag checks fail
because $\mathit{TreeEq}(\mathit{Pair}(\mathit{Pair}(\ldots), \ldots),
\mathit{reify}(\mathit{natCode}\;j)) = \mathit{Pair}(O,O)$ for all~$j$.
The dispatch falls through to $\mathit{sndArg2}$, which returns the
recursive results: $\mathit{Pair}(\mathit{thFunTFor}(a),
\mathit{thFunTFor}(b))$.
\end{definition}

The passthrough is proved in $\mathit{Deriv}$ by
$\mathit{ndDispatchPairMiss}$ (from \texttt{PairPassthroughAll.agda}).

\subsection{Experiment A: axRefl (tag~17)}

\begin{theorem}[nelsonAxRefl]
For all $x$:
\[
  \mathit{thFunTFor}(\mathit{Pair}(\mathit{tag}_{17},
  \mathit{Pair}(x, O))) = \mathit{Pair}(x, x)
\]
is derivable in $\mathit{Deriv}$.
\end{theorem}

This says: the proof checker correctly handles reflexivity encodings.
The output $\mathit{Pair}(x,x)$ is the code of $x = x$. The proof
dispatches through 17~branch misses, hits tag~17, and reduces the
case-17 extractor $\mathit{origA}$ to~$x$. Compile time: $0.12$s.

\subsection{Experiment B: ruleTrans (tag~19)}

\begin{theorem}[nelsonRuleTrans]
For all $\mathit{sp}_1$, $\mathit{sp}_2$:
\[
  \mathit{thFunTFor}(\mathit{Pair}(\mathit{tag}_{19},
  \mathit{Pair}(\mathit{sp}_1, \mathit{sp}_2)))
  = \mathit{Pair}(\mathit{Fst}(\mathit{Fst}(\mathit{thFunTFor}(
  \mathit{Pair}(\mathit{sp}_1, \mathit{sp}_2)))),\;
  \mathit{Snd}(\mathit{Snd}(\mathit{thFunTFor}(
  \mathit{Pair}(\mathit{sp}_1, \mathit{sp}_2)))))
\]
is derivable in $\mathit{Deriv}$.
\end{theorem}

This is the raw result: the proof checker extracts from the recursive
evaluation of the sub-proof pair.

\begin{theorem}[nelsonRuleTransFull]
When $\mathit{sp}_1 = \mathit{Pair}(\mathit{Pair}(a_1,a_2), b_1)$ and
$\mathit{sp}_2 = \mathit{Pair}(\mathit{Pair}(c_1,c_2), d_1)$
(Pair-tagged, ensuring passthrough):
\[
  \mathit{thFunTFor}(\mathit{Pair}(\mathit{tag}_{19},
  \mathit{Pair}(\mathit{sp}_1, \mathit{sp}_2)))
  = \mathit{Pair}(\mathit{Fst}(\mathit{thFunTFor}(\mathit{sp}_1)),\;
  \mathit{Snd}(\mathit{thFunTFor}(\mathit{sp}_2)))
\]
is derivable in $\mathit{Deriv}$.
\end{theorem}

This is the full compositional result: the proof checker on a
$\mathit{ruleTrans}$ encoding returns the left side of the first
sub-proof's equation paired with the right side of the second.
The proof composes tag dispatch, the passthrough lemma, and extractor
reduction. Compile time: $0.13$s.

\subsection{Experiment C: ruleInst (tag~23)}

\begin{theorem}[nelsonRuleInst]
For Pair-tagged $\mathit{tC}$ and $\mathit{sp}$:
\[
  \mathit{thFunTFor}(\mathit{Pair}(\mathit{tag}_{23},
  \mathit{Pair}(\mathit{Pair}(\mathit{tC}, \mathit{xC}), \mathit{sp})))
  = \mathit{Pair}(\mathit{substTFor}(\mathit{Fst}(\mathit{thFunTFor}(\mathit{sp}))),\;
  \mathit{substTFor}(\mathit{Snd}(\mathit{thFunTFor}(\mathit{sp}))))
\]
is derivable in $\mathit{Deriv}$.
\end{theorem}

This is the substitution case: the proof checker applies
$\mathit{substTFor}$ (the object-level coded substitution, with free
variables $v_{11}$, $v_{12}$) to both sides of the sub-proof's equation.
This is the case where the self-referential coding apparatus enters
the proof checker. Compile time: $0.14$s.

\subsection{Experiment D: ruleF / Schema~F (tag~24)}

\begin{theorem}[nelsonRuleF]
For all $\mathit{fC}$, $\mathit{gC}$, $\mathit{data}$:
\[
  \mathit{thFunTFor}(\mathit{Pair}(\mathit{tag}_{24},
  \mathit{Pair}(\mathit{Pair}(\mathit{fC}, \mathit{gC}), \mathit{data})))
  = \mathit{Pair}(\mathit{Pair}(\mathit{tag}_{\mathit{ap1}},
  \mathit{Pair}(\mathit{fC}, v_0)),\;
  \mathit{Pair}(\mathit{tag}_{\mathit{ap1}},
  \mathit{Pair}(\mathit{gC}, v_0)))
\]
is derivable in $\mathit{Deriv}$, where $v_0 = \mathit{reify}(\mathit{code}(\mathit{var}\;0))$.
\end{theorem}

The proof checker on a Schema~F encoding produces the coded equation
$f(x) = g(x)$. This is the system verifying its own induction principle.
Compile time: $0.12$s.

\subsection{Schema~F Internal Application}

\begin{theorem}[passFIsId]
The passthrough function $\mathit{Rec}\;O\;\mathit{sndArg2}$ equals the
identity~$I$ on all trees:
\[
  \mathit{Rec}(O, \mathit{sndArg2})(t) = t
\]
is derivable in $\mathit{Deriv}$ for all~$t$, via Schema~F.
\end{theorem}

This verifies that Schema~F can be used internally to prove properties
of recursive functions. The proof defines both $\mathit{Rec}\;O\;\mathit{sndArg2}$
and~$I$ as satisfying the same recursion (base $O$, step $\mathit{sndArg2}$),
then applies $\mathit{ruleF}$.

\section{Analysis}

\subsection{What Succeeds}

All four structurally distinct case types of the proof checker are
internally verifiable:

\begin{center}
\begin{tabular}{llll}
\hline
\textbf{Experiment} & \textbf{Case} & \textbf{Structure} & \textbf{Time} \\
\hline
A: axRefl & tag 17 & Base (ignores recs) & 0.12s \\
B+: ruleTrans & tag 19 & Recursive + passthrough & 0.13s \\
C: ruleInst & tag 23 & Substitution (substTFor) & 0.14s \\
D: ruleF & tag 24 & Induction principle & 0.12s \\
\hline
\end{tabular}
\end{center}

The remaining 22 cases are structural variations of these four patterns.
The instant compile times indicate that these are native equational
calculations for the system, not edge cases.

\subsection{Assembly Obstacles}

There are two levels of assembly.

\textbf{Assembly~1} (first-order): combine the 26 branch lemmas into a
single internal theorem stating that $\mathit{thFunTFor}$ obeys its
decoding specification on all trees. This would require packaging the
26-case step function as a single $\mathit{Fun2}$ combinator and using
Schema~F. The obstacle is size (enormous but finite), not conceptual.

\textbf{Assembly~2} (soundness): go from the decoding theorem to an
internal consistency statement. This requires bridging the gap between
``$\mathit{thFunTFor}$ computes this equation'' (first-order) and
``this equation is derivable'' (involves $\mathit{Deriv}$, which is a
meta-level type). This is the likely real failure point for Nelson's
program in this system.

\subsection{Nelson's Response}

Nelson would collapse the gap by arguing that the encoding tree \emph{is}
the derivation witness. Soundness becomes: ``$\mathit{thFunTFor}$ correctly
decodes every tree,'' which is first-order. The experiments show this
decoding is internally verifiable case-by-case. Whether it can be
assembled into a global theorem via Schema~F---and whether that global
theorem triggers the G\"odel~II contradiction---remains the decisive
open question.

\section{Supplementary: Substitution Analysis}

\subsection{The 108s instForm2 Computation}

The only heavy computation in the G\"odel~II formalization is
\texttt{instForm2} in \texttt{SubstTForPrecomp.agda} (108 seconds,
cached). This proves by \texttt{refl} that double substitution
($v_{11}$, $v_{12}$) on an equation containing $\mathit{cstf}$ and
$\mathit{substTFor}$ yields the expected form. The cost is
materializing $\mathit{reify}(\mathit{crTC})$ (${\sim}400$K Term nodes).

\subsection{Schematic Lemmas (Instant)}

We prove schematic substitution lemmas that avoid materializing any
large tree:

\begin{itemize}
\item $\mathit{substReify}$: $\mathit{subst}\;x\;r\;(\mathit{reify}\;t)
  = \mathit{reify}\;t$ for all $t$ (by tree induction).
\item $\mathit{substTForClose}$: double substitution on $\mathit{substTFor}$
  gives $\mathit{closedSubstTFor}$ (schematic, uses $\mathit{substReify}$
  only on the small $\mathit{tgtT}$, not on $\mathit{crTC}$).
\item $\mathit{thFunTForSubst}$: double substitution on $\mathit{thFunTFor}$
  replaces the embedded $\mathit{substTFor}$ (in case~23) with
  $\mathit{closedSubstTFor}$, propagated through the 26-branch chain.
\item $\mathit{instFormGen}$: general bridge lemma for the ruleInst step
  of G\"odel~II, combining the above.
\end{itemize}

All compile in under $0.2$~seconds. The 108s computation has no
counterpart in Rose's or Guard's mathematical proofs: they never
materialize a specific numeral, but reason schematically about the
coding functions. The schematic lemmas restore this correspondence.

\subsection{Key Insight}

The 108s is not caused by substitution complexity. It is caused by
$\mathit{cstf}$ appearing in the \emph{type} of $\mathit{instForm2}$,
which forces Agda to unfold $\mathit{cstf} =
\mathit{closedSubstTFor}(\mathit{reify}\;\mathit{crTC}, \ldots)$ and
thereby materialize the ${\sim}400$K-node term. Restructuring the
G\"odel~II proof to perform $\mathit{ruleInst}$ before introducing
$\mathit{cstf}$ (matching Rose's proof order) would eliminate this
computation entirely.

\section{File Inventory}

\begin{center}
\begin{tabular}{lrl}
\hline
\textbf{File} & \textbf{Lines} & \textbf{Role} \\
\hline
DerivSize.agda & 85 & derivSize + Schema F size = 113/125 \\
SubstReify.agda & 120 & substReify + substClosedSTF \\
SubstTForClose.agda & 45 & substTForClose (v12-first) \\
ThFunTForSubst.agda & 210 & thFunTForSubst + thFunTForClosed \\
InstForm.agda & 50 & instFormGen (general bridge) \\
NelsonAxRefl.agda & 220 & Experiment A: axRefl \\
NelsonRuleTrans.agda & 175 & Experiment B: ruleTrans (raw) \\
NelsonRuleTransFull.agda & 160 & Experiment B+: ruleTrans (full) \\
NelsonRuleInst.agda & 170 & Experiment C: ruleInst \\
NelsonRuleF.agda & 165 & Experiment D: ruleF \\
NelsonSchemaF.agda & 65 & Schema F: passF = I \\
\hline
\end{tabular}
\end{center}

Total: approximately 1465 lines of new Agda code. All files compile in
under 0.2 seconds (excluding cached imports).

\begin{thebibliography}{9}
\bibitem{rose1961} H.\,E.\ Rose,
\emph{On the Consistency and Undecidability of Recursive Arithmetic},
Zeitschr.\ f.\ math.\ Logik und Grundlagen d.\ Math., Bd.~7,
S.~124--135 (1961).

\bibitem{guard1963} J.\,R.\ Guard,
\emph{Lecture Notes on Recursive Arithmetic},
Princeton University, Mathematics Department (1963).

\bibitem{nelson} E.\ Nelson,
\emph{Predicative Arithmetic},
Mathematical Notes~32, Princeton University Press (1986).
\end{thebibliography}

\end{document}
